\section{Discussion}
\label{sec:discussion}

The pattern in Table~\ref{tab:summary} reveals an under-appreciated property of strong fine-tuned visual classifiers on biological data: \emph{information-theoretic orthogonality is not the same as statistical orthogonality} once the visual model has been trained on data where the auxiliary signal is naturally correlated with image content. cRT was structurally orthogonal but introduced wrong evidence and dragged the anchor down by $0.036$; genus reweighting was statistically dependent on the BioCLIP 2.5 posterior itself and moved the score by $-0.006$, close to the noise floor; the seasonal phenology prior --- the most genuinely orthogonal source we tested, since date is not in the image at all --- moved the score by only $-0.005$, again close to the noise floor, because the seasonal structure was already implicit in the leaf and flower phenotypes seen during training.

The practical implication is that real headroom on this benchmark requires structurally stronger \emph{visual} models (end-to-end fine-tuning that preserves the prior, multi-crop test-time augmentation, test-time training, synthetic mosaic training distributions) rather than post hoc reweighting with auxiliary priors. A natural next experiment, blocked here by the lack of test site coordinates, is a climate-envelope variant: per-species 12-D bioclim Gaussians fit from training image lat/lons, scored against WorldClim variables at each test site --- this would give a prior conditioned on \emph{location} as well as date, which is the one orthogonal axis that visual features cannot recover from the image alone.

We also observed a striking decoupling between validation top-5 accuracy on the held-out 10\% split of the single-plant manifest and Macro F1 on the multi-species quadrat leaderboard. Across our twelve strongest runs the two metrics correlate at only $r \approx 0.4$ (Figure~\ref{fig:val_vs_kaggle}), and within the unfrozen-blocks ablation of Table~\ref{tab:blocks} they invert: validation top-5 accuracy rises monotonically with unfrozen capacity while leaderboard Macro F1 peaks in the middle and collapses at full fine-tune. We read this as direct evidence that the bottleneck on this benchmark is the single-plant to multi-species distribution shift, not model capacity on the single-plant distribution itself.

\begin{figure}[!htbp]
  \centering
  \input{figures/val_vs_kaggle}
  \caption{Validation top-5 accuracy on the held-out single-plant split versus public Macro F1 on the multi-species quadrat leaderboard, across our twelve strongest runs. The headline anchor ($\star$, i002 last-four-blocks 224+336+LA ensemble) does not sit at the maximum validation point; the i003 capped checkpoints (green squares) and the 015 PC24+iNat mix (teal) reach higher validation accuracy but lower leaderboard score. The dashed line indicates the weak monotone trend ($r \approx 0.4$).}
  \label{fig:val_vs_kaggle}
\end{figure}
